\documentclass[../../main]{subfiles}

\usepackage{booktabs}
%\usepackage{array}
%\usepackage{parskip}
%\usepackage{setspace}
%\usepackage{enumitem}

\newenvironment{dsttab}[1]{%
  \begingroup\setlength\tabcolsep{3pt}\small
  \setlength{\extrarowheight}{2pt}%
  \begin{longtable}{#1}%
}{%
  \end{longtable}\endgroup
}

\title{A European Digital Services Tax: The Case for an EU Instrument}
\date{July 2026}
\author{Samuel Guerin}

\begin{document}

\maketitle
\thispagestyle{empty}

Profit shifting and tax optimisation specific to the digital sector enable digital companies to face a lower effective corporate income tax rate than companies in other sectors.
In the EU, this gap is substantial and raises questions about the fairness of our tax system.
The uncollected tax revenues also represent a significant amount of money that escapes Member States' budgets.
Taxing those profits at source requires international consensus, but six years of OECD negotiations have failed to produce one.
Several EU member states have turned instead to taxing digital revenues, a second-best instrument that is nonetheless administrable and politically tractable.

This report makes the case for replacing that fragmented landscape with a single EU-level instrument structured as an Own Resource: a 3\% levy on digital advertising, intermediation, and user data revenues.
The instrument acts as a proxy for the corporate income tax revenues that are not collected from digital companies, yielding approximately EUR~6\,bn annually by 2028.
In the same vein as the national Digital Service Taxes that it will replace, it is a temporary measure enacted to foster the international process towards profit taxation.
It also mitigates individual Member States' exposure to US trade pressure by leveraging a single Union-level stance.

The report also proposes another option to contribute to the repayment of NextGenerationEU, following a proposal by \cite{acemoglu_johnson_2024_ad_tax}: a Pigouvian levy targeted at digital advertising.
This instrument targets the externalities generated by the attention-capture business model in this sector. It is designed to redirect economic activity and technological development away from activities associated with these documented externalities, including harms to individual mental health and threats to democratic discourse.
We propose an initial tax rate of 15\%, which would raise EUR~12--14bn annually by 2028, as a first step towards the 50\% tax rate proposed by Acemoglu and Johnson.

% ------------------------------------------------------
\section{The problem: the corporate income tax was not built for the digital economy}
\label{sec:problem}


The digitalisation of the economy challenges our existing tax system.
The system is based on the principle that companies are taxable where they maintain a sufficiently substantial physical presence.
That doctrine was designed in the 1920s for an industrial economy and has not been revised in any structural sense since.
%Digital services, being less anchored to any physical location, allow companies to shift profits to low-tax jurisdictions through legal but economically opaque intragroup arrangements.
Technology companies can generate billions in advertising revenues, intermediation fees, and user data monetisation from users in a country and shift profits via royalty payments and intercompany transfer prices to entities holding intellectual property in another country.
Not only does this amount to a transfer of corporate income tax revenues away from the places where value is created, but at a global scale, the fiscal consequence is measurable.
The European Commission estimated in its 2018 impact assessment that digital companies operating in the EU face an effective corporate income tax rate of 9.5\%, against 23.2\% for traditional companies of comparable scale~\citep{swd_2018_81}.\footnote{Figures based on modelled effective average tax rates (ZEW/PwC 2016 methodology). Pillar Two (Directive 2022/2523/EU, FY2024) establishes a 15\,\% floor but does not close the structural gap; no updated EU-wide estimate has been published.}

The debate on digital taxation in the European Union emerged from this growing fiscal mismatch.
Member states raised concerns about tax fairness, competitive neutrality, and revenue losses for governments.
At the same time, discussions on new EU own resources brought additional attention to the issue.
In March 2018, the European Commission proposed two legislative initiatives: a significant digital presence rule to establish a digital permanent establishment, and a Digital Services Tax (DST) as a temporary measure pending an international agreement.
The proposed DST would have applied a 3\% tax to revenues from online advertising, digital intermediation services, and the sale of user-generated data, targeting only multinational groups with annual worldwide revenues above EUR 750 million and EU digital revenues above EUR 50 million.
Estimated annual revenues stood between EUR 4.7 and 6 billion.
Although the proposal failed to secure unanimous support in the Council, it inspired several national DSTs. 

These concerns were not confined to Europe, and calls grew for a reform of the international tax system.
These calls led the G20 to mandate the OECD to develop solutions to address base erosion and profit shifting (BEPS). 
Following the completion of the OECD/G20 BEPS Project, negotiations continued within the OECD/G20 Inclusive Framework, which brings together more than 140 jurisdictions to further develop international tax rules. 
These negotiations culminated in the two-pillar solution. 
Pillar One aims, in particular, to reallocate taxing rights to market jurisdictions.
It applies to the largest multinational enterprises, with global revenues exceeding USD 20 billion and profitability above 10\%. 
Under its rules, 25\% of residual profits (profits exceeding the 10\% profitability threshold) are allocated to market jurisdictions in proportion to the revenues generated in each jurisdiction.

%The OECD responded with the two-pillar BEPS~2.0 framework. Pillar~1 (Amount~A) would have reallocated 25\% of the residual profits of the largest multinationals (global revenues above USD~20 billion, margins above 10\%) to the market jurisdictions where users are located. In exchange, countries maintaining a digital services tax would commit to withdrawing it upon implementation. Pillar~2 would establish a global minimum effective tax rate of 15\% through the Income Inclusion Rule and the Undertaxed Profits Rule.

The Multilateral Convention was designed to implement this reallocation of taxing rights, superseding bilateral tax treaties between signatory jurisdictions.
Although the text has been finalised, ratification difficulties stalled the project.
The United States withdrew from the negotiations in January 2025. 
The co-chair statement of the Inclusive Framework acknowledged that the MLC text itself is stable but that unresolved issues prevent conclusion. 
No timeline exists. 

A parallel process is under way at the United Nations, initiated by a large group of developing countries considering that the OECD framework did not adequately represent their interests.
The UN General Assembly adopted the Terms of Reference for a Framework Convention on International Tax Cooperation (UNFCITC) in 2024, opening an intergovernmental negotiation expected to run until 2027.
The convention's ambition goes well beyond digital taxation: its goal is to establish a new architecture for the global distribution of taxing rights, with an emphasis on the source-based taxation principle.
Protocol I of the convention addresses the taxation of income from cross-border services. 
Digitalised services are included, with a broader scope than that of the European DSTs already implemented and our proposal.
As of mid-2026, negotiations show no signs of concluding. % AF: si tu veux j'ai les draft treaty qui sont en train d'être négociés

European DSTs emerged as national responses to the slow pace of OECD/G20 negotiations, and their proliferation subsequently increased US pressure to conclude Pillar One.
Participating jurisdictions of Pillar One agreed to withdraw their DSTs upon implementation of the new agreement.
The moratorium on new DSTs, maintained by OECD members as a confidence-building measure during the Pillar One negotiations, expired without renewal at the end of 2023. 
Countries are therefore no longer under a formal commitment to refrain from introducing or expanding DSTs. 
As those incentives weakened, the United States continued to pressure trading partners to abandon DSTs through the threat of trade and economic measures.

The debate on digital taxation is embedded in a complex international framework, and a European DST is designed to address concerns that extend well beyond Europe's borders.
It will change the relationship with major trading partners, and pressure should be expected.
But the instrument also pushes for more global cooperation: by acting as a single bloc, the EU can demand a level of international alignment that no member state could claim alone.

% -------------------------------------------------------
\section{National digital services taxes: the current landscape in the European Union}
\label{sec:landscape}

The European Council's failure to reach unanimity on the 2018 proposal by the Commission left member states with a choice: wait for an international agreement that remained uncertain, or act unilaterally.
Several chose to act.
France enacted the \emph{taxe sur les services numériques} in 2019, triggering a wave of similar instruments across Europe.
Seven years later, six jurisdictions have enacted DSTs, each with its own rate, scope, and threshold structure (Hungary's rate has been suspended at 0\% since April 2026).

France, Italy, and Spain drew directly on the Commission's template, adopting the same three-category tax base (advertising, intermediation, and user data) at the same 3 per cent rate on revenues.
Austria chose a narrower instrument: online advertising only, at 5 per cent.
Hungary stands apart: its advertising tax dates to 2014, predates the DST wave, and originated as a progressive media levy with motivations that were as much political as fiscal.
In parallel, the United Kingdom, acting outside the EU framework after Brexit, introduced a DST at 2 per cent in 2020 on advertising and intermediation revenues, calibrated in part to limit its exposure to the United States Trade Representative (USTR) Section~301 process.


% --- Experimental table: factored scope + split threshold ---
\begin{dsttab}{@{}p{2.4cm}p{1.2cm}cccccp{1.8cm}p{2.7cm}@{}}
\toprule
\textbf{Country} & \textbf{Rate} & \multicolumn{4}{c}{\textbf{Scope}} & \multicolumn{2}{c}{\textbf{Threshold (EUR\,M)}} & \textbf{Revenue} \\
\cmidrule(lr){3-6}\cmidrule(lr){7-8}
 & & \textbf{Ad.} & \textbf{Int.} & \textbf{Data} & \textbf{Stream.} & \textbf{Global} & \textbf{Domestic} & \textbf{(EUR\,M)} \\
\midrule
\endfirsthead
\multicolumn{9}{l}{\textit{(continued)}} \\
\toprule
\textbf{Country} & \textbf{Rate} & \multicolumn{4}{c}{\textbf{Scope}} & \multicolumn{2}{c}{\textbf{Threshold (EUR\,M)}} & \textbf{Revenue (EUR\,M)} \\
\cmidrule(lr){3-6}\cmidrule(lr){7-8}
 & & \textbf{Ad.} & \textbf{Int.} & \textbf{Data} & \textbf{Stream.} & \textbf{Global} & \textbf{Domestic} & \\
\midrule
\endhead
\midrule
\multicolumn{9}{r}{\textit{(continued on next page)}}
\endfoot
\endlastfoot
France (TSN)  & 6\%                        & $\checkmark$ & $\checkmark$ & $\checkmark$ &              & 750 & 25       & 756\,(2024) \\
Italy (ISD)   & 3\%                        & $\checkmark$ & $\checkmark$ & $\checkmark$ &              & 750 & ---      & 637\,(2025) \\
Spain (IDSD)  & 3\%                        & $\checkmark$ & $\checkmark$ & $\checkmark$ &              & 750 & 3        & 375\,(2024) \\
Austria       & 5\%                        & $\checkmark$ &              &              &              & 750 & 25       & 124\,(2024) \\
Hungary       & 7.5\%\textsuperscript{a}   & $\checkmark$ &              &              &              & --- & 100      & ---         \\
%Turkey       & 5\% (2026)                 & $\checkmark$ & $\checkmark$ &              & $\checkmark$ & 750 & TRY\,20  & ---         \\
\hline
UK (DST)      & 2\%                        & $\checkmark$ & $\checkmark$ &              &              & GBP\,500 & GBP\,25 & GBP\,944\,(2026) \\
\bottomrule
\multicolumn{9}{p{15cm}}{\footnotesize Sources: \citet{ceps_april2025,corte_dei_conti_rendiconto_2025,aeat_informe_anual_2024,tax_foundation_dst_tracker_2026,rtr_digitalsteuer_2025}.
\textsuperscript{a} statutory 7.5\%; suspended to 0\%, April 2026. Non-EUR domestic thresholds shown with currency symbol.} \\
\caption{Digital services taxes in force, mid-2026 (factored)}\label{tab:dst2}
\end{dsttab}

National DST revenues have proved stable and growing in the European Union.
The UK's \emph{Digital Services Tax} revenues rose from GBP 358 million in 2020--21 to GBP 944 million in 2025--26, a five-year CAGR % AF: define CAGR
of 21.4\%, driven purely by the expansion of the digital economy as the tax rate has not changed since implementation.
Italy's ISD grew from EUR 298 million (2022) to EUR 637 million (2025), a 113\% increase over three years, accelerated by the January 2025 removal of its domestic revenue threshold.
France doubled its rate from 3\% to 6\% in the 2026 finance law~\citep{legifrance_plf2026,assemblee_nationale_plf2026_amdt655}, motivated by fiscal need and mounting frustration with Pillar~1's collapse.
Those cases show that DST receipts draw on two sources: the structural expansion of the digital economy, and deliberate adjustments to the instruments themselves.
Those adjustments encode a political position: a country raises its DST to build leverage and force the negotiation forward; it reduces or withdraws when bilateral pressure makes the instrument too costly to hold, or when a deal becomes credible enough to justify stepping back.
These instruments are not designed to last.
They are designed to move.

The same logic operates beyond Europe but in an opposite direction.
Turkey reduced its rate from 7.5\% to 5\% in January 2026 and has scheduled a further cut to 2.5\% from 2027~\citep{tax_foundation_dst_tracker_2026}, using the rate trajectory as a diplomatic signal of readiness to step back.
Canada implemented a 3\% DST in June 2024 and rescinded it twelve months later under direct bilateral pressure from the United States~\citep{canada_finance_rescission_jun2025}.
India abolished both of its digital levies in 2024 and 2025~\citep{india_briefing_el2_abolition,india_briefing_el1_abolition}, as a unilateral bet that OECD implementation would follow.

National experience shows that the instrument is administratively tractable: applied to a small number of large multinationals on a gross revenue base, it has proved collectible in practice, and high thresholds have consistently excluded SMEs from its scope.

The European Union framework is dynamic but fragmented.
The majority of EU member states levy no digital services tax and extract nothing from digital revenues generated in their territories, even as those revenues grow.
Among those that do, each instrument sets its own thresholds, tax base, rate, and enforcement mechanism.
Rates vary and definitions of ``digital advertising'' or ``intermediation service'' may differ across jurisdictions.
For example, the same app distribution platform is a taxable intermediary in France and a non-taxable digital content supplier in Italy and Spain.  
The resulting divergence creates legal uncertainty and compliance costs for multinationals operating across member states.
It also creates structuring incentives to escape those national instruments.
Finally, individual member states implementing a DST have faced US Section~301 bilateral pressure without the collective leverage of the Union.




Beyond the instruments in force, Belgium's coalition government has committed to a 3\% DST by 2027, contingent on the absence of EU or OECD consensus~\citep{ceps_april2025}; Poland has a 3\% proposal at legislative review; Slovakia published a formal proposal in August 2025; the Czech Republic has a 7\% proposal stalled in parliament.
The question for the next few years is not whether EU member states will tax digital services.
They will.
The question is whether they do so in a fragmented national patchwork or through a coordinated EU instrument.

\section{A European digital services tax}
The proposed instrument is structured as an EU own resource under Article~311 TFEU: revenues feed the EU budget directly rather than national treasuries, and national tax administrations act as collection agents on the Union's behalf.
The design follows the 2018 Commission proposal~\citep{eu_commission_proposal} with all revenues directed to the EU budget.
Three categories of digital services are covered: targeted advertising, intermediation, and the transmission of user data.
A flat rate of 3 per cent applies to gross revenues before any deduction for costs, intragroup charges, or IP royalties.
The three categories share a common characteristic: the users themselves are the productive input.
Advertising platforms sell attention; marketplaces sell access to a buyer and seller base; data services sell records of user behaviour.
In each case, the platform captures and monetises value that users generate collectively.

\paragraph{Targeted advertising.}
Advertising is targeted if the selection of which ad appears to a given user relies on data collected from that user's prior activity.
The charge falls on the entity providing the placement service: where an intermediary places ads on third-party publisher interfaces, the intermediary's revenues are taxable and the publisher incurs no second charge on the same amount.

\paragraph{Intermediation services.}
The taxable service is the operation of a multi-sided digital interface through which users locate and interact with other users.
The scope consists of matching platforms and marketplace platforms.
The taxable revenues are fees charged for access, promoted visibility or data-based matching .
For transactions, only the service fee is charged, not the value of that transaction.

\paragraph{User data transmission.}
The taxable service is the sale of data collected from users' activity on digital interfaces to a third party.
The perimeter is narrow: only data collected from users' digital activity, transmitted to a third party, and in exchange for payment.
Data used internally, shared without charge, or not originating from users' activity are all outside scope.
In practice this is the smallest of the three categories.
The dominant platforms primarily monetise behavioural data through advertising, already covered by the first category; an independent charge arises mainly for data broker activity where the aggregator and the advertising buyer are separate entities.

\paragraph{Exclusions.}
The rationale is that value for content providers derives from the content itself, not from user co-creation; the multi-sided network dynamic that justifies the levy is absent.
Excluded from scope is digital content supplied to users through a digital interface, including video, audio, games, news, and text, as well as cloud computing and software-as-a-service.

\paragraph{Nexus.}
Taxing rights attach to the EU for revenues derived from users located within the EU at the moment of the taxable interaction, identified by the IP address of the device.
No permanent establishment or physical presence in any member state is required.
The user-location principle has not been successfully challenged before any EU member state court since national DSTs entered into force in 2019, establishing a robust precedent for adoption at EU level.

\paragraph{Thresholds.}
The tax applies only to groups exceeding EUR~750~million in total worldwide revenues from all sources and EUR~50~million in DST taxable revenues within the EU.
Both tests apply at consolidated group level, preventing threshold avoidance through corporate fragmentation.

\paragraph{Sunset clause.}
The instrument includes an automatic repeal on the date of entry into force of any international agreement on digital economy taxation, making it explicitly transitional.


%Transfer-pricing optimisation that reduces corporate income tax by routing profits through low-tax IP holding companies does not affect DST liability: advertising impressions in Germany and marketplace transactions in France remain taxable regardless of where the revenue is invoiced.
%The proposed rate is 3 per cent.
%At this rate, the three-category EU-27 base yields EUR~4.3\,bn per year (Table~\ref{tab:compare}).

\section{Revenue estimates}
\label{sec:revenue_estimates}

We estimate the revenues that could be generated by such a DST.
For analytical purposes, we allocate these revenues across Member States.
This requires solving a geographical allocation problem that no public source addresses directly: revenues from digital services are booked by legal entity, typically in the jurisdiction where the operating subsidiary is registered, rather than where users are located.
%Three estimation strategies appear in the literature.

\paragraph{Top-down market aggregate (M1).}
Rather than estimating each country separately, this method starts from a single Commission-estimated EU-wide aggregate and distributes it across member states using each country's share of the digital advertising market as a proxy:

\begin{equation}
  R_i = \tau \cdot s_i \cdot M
  \label{eq:topdown}
\end{equation}
The EU-wide three-category revenue base $M$ is distributed among countries according to their respective shares of the European digital advertising market, based on IAB Europe's 2024 national data~\citep{iab_adex_2024}: $s_i = A_i / A_{\text{EU-27}}$ where $A_i$ is what advertisers spent in a market $i$.
In Staff Working Document SWD(2018)~81~\citep{swd_2018_81} the base revenue is obtained by aggregating firm-level revenues across the three in-scope categories for EU-28 in 2019, yielding EUR~4.7\,bn.
The UK accounts for roughly 30\% of that aggregate (SWD footnote~71); removing it gives an EU-27 base of EUR~3.3\,bn for 2019.
To update to 2024, we scale by the IAB EU-27 advertising market growth 2022--2024 (factor 1.363)~\citep{iab_adex_2024}, giving $M = \text{EUR}~4{,}484\,\text{M}$.
Country shares $s_i$ use IAB 2024 national advertising data, so the allocation reflects current market geography while the aggregate follows the SWD methodology.\footnote{Two new disclosure regimes improve the outlook. The EU Public CbCR Directive (2021/2101/EU) requires groups with revenues above EUR~750\,million to publish revenues by EU member state from FY2024 onward. The GloBE Information Return, filed for the same threshold from 2026, contains jurisdiction-level revenue data for all in-scope groups, though it remains confidential between administrations. Together, these regimes will make a bottom-up estimation of~$M$ tractable for future work.}
The principal weakness remains that market data record where revenues are invoiced rather than where users are located, concentrating estimated revenues in countries with large subsidiary operations (Ireland, the Netherlands, Sweden) and understating them in large consumer markets such as Poland and Romania.

% AF: I wouldn't call them M1, M2 but "top-down" or "invoices", "bottom-up", etc.
\paragraph{Bottom-up calibration from DST returns (M2).}

We can also use the same data, together with revenues from already implemented DSTs, to provide a bottom-up estimate of the revenues that could be generated by an EU-wide DST.
We assume that intermediation fees and user data transmission revenues are proportional to advertising ones:
\begin{equation}
  R_i^{\mathrm{est}} = \tau \cdot \bar{k} \cdot A_i
  \label{eq:M2}
\end{equation}
France and Spain both operated a three-category DST at 3\% through end-2025; their confirmed 2024 revenues at that rate are the natural anchors for $\bar{k}$.\footnote{France doubled its rate to 6\% in the 2026 finance law; the calibration uses 2024 data collected at 3\%.}
With $A_{\mathrm{FR}} = \text{EUR}~11{,}209\,\text{M}$ and $A_{\mathrm{ES}} = \text{EUR}~5{,}501\,\text{M}$:
\begin{equation}
  k_i = \frac{R_i^{\mathrm{obs}} / \tau_i}{A_i}, \qquad
  k_{\mathrm{FR}}  = 2.25, \qquad
  k_{\mathrm{ES}}  = 2.27
  %k_{\mathrm{FR}} = \frac{756/0.03}{11{,}209} = 2.25, \qquad
  %k_{\mathrm{ES}} = \frac{375/0.03}{5{,}501} = 2.27
  \label{eq:calib}
\end{equation}

The mean $\bar{k} = 2.26$: the three-category base is approximately 2.3 times the advertising market alone.
Applied to the IAB EU-27 total, this gives an aggregate of EUR~4{,}298\,M.
Austria (5\% advertising-only, EUR~124\,M) gives $k_{\mathrm{AT}} = 1.13$, consistent with a single-category instrument, the residual gap above unity reflecting the broader legal definition of taxable advertising revenues relative to IAB survey coverage; Italy's 2025 figure implies $k_{\mathrm{IT}} = 3.30$, inflated by the January 2025 threshold removal.
France and Spain remain the preferred anchors.
$k_{\mathrm{FR}} = 2.25$ and $k_{\mathrm{ES}} = 2.27$ differ by less than 1\,\%, making the EU-27 aggregate insensitive to the choice of anchor.

\paragraph{Consumer-location reallocation (M3).}
M3 uses the same EU-27 aggregate as M2 but applies a different allocation key.
Where M2 distributes by IAB national advertising shares (booking geography, headquarters-weighted), M3 reallocates by demand-side weights: each country's share of online buyers in the EU-27, from Eurostat's survey~\citep{eurostat_isoc_ec_ib20} (individuals aged 16--74 ordering online in the last three months, multiplied by national population).
The key matches the intermediation nexus directly; for the advertising category, the correct variable would be ad impressions by device location, which is not publicly available.
Online buyers remain the best public proxy for the combined three-category base.

Table~\ref{tab:compare} presents member-state estimates under all three methods.
M2 and M3 both yield an EU-27 total of EUR~4.3\,bn, grounded in observed national DST revenues; they differ in how that total is distributed across member states.

\begin{dsttab}{@{}p{3.8cm}rrrrr@{}}
\toprule
\textbf{Member state} & \textbf{M1} & \textbf{M2} & \textbf{M3} & \textbf{Confirmed} & \textbf{M3 (2028)} \\
 & \textit{\footnotesize SWD/2024} & \textit{\footnotesize calibr./2024} & \textit{\footnotesize Eurostat/2024} & \textit{\footnotesize admin.} & \textit{\footnotesize proj.} \\
\midrule
\endfirsthead
\multicolumn{6}{l}{\textit{(continued)}} \\
\toprule
\textbf{Member state} & \textbf{M1} & \textbf{M2} & \textbf{M3} & \textbf{Confirmed} & \textbf{M3 (2028)} \\
\midrule
\endhead
\midrule
\multicolumn{6}{r}{\textit{(continued on next page)}}
\endfoot
\endlastfoot
Austria              & 155       & 148       & 104       & 124$^{a}$       & 141        \\
Belgium              & 86        & 82        & 147       & ---       & 200       \\
Bulgaria             & 8         & 8         & 28        & ---       & 38        \\
Croatia              & 12        & 12        & 34        & ---       & 46        \\
Cyprus         & ---$^{b}$       & ---       & 12        & ---       & 16        \\
Czech Republic       & 123       & 118       & 130       & ---       & 177       \\
Denmark              & 121       & 116       & 78        & ---       & 106       \\
Estonia              & 10        & 10        & 17        & ---       & 23        \\
Finland              & 63        & 60        & 71        & ---       & 97        \\
France         & 793       & 760$^{c}$       & 792       & 756       & 1{,}078   \\
Germany              & 1{,}265   & 1{,}213   & 715       & ---       & 973       \\
Greece               & 26        & 25        & 67        & ---       & 91        \\
Hungary              & 53        & 51        & 102       & ---       & 139       \\
Ireland              & 75        & 72        & 71        & ---       & 97        \\
Italy          & 417       & 400       & 341       & 637$^{d}$       & 464       \\
Latvia               & 9         & 9         & 20        & ---       & 27        \\
Lithuania            & 11        & 11        & 30        & ---       & 41        \\
Luxembourg     & ---$^{b}$       & ---       & 7         & ---       & 10        \\
Malta          & ---$^{b}$       & ---       & 5         & ---       & 7         \\
Netherlands          & 306       & 293       & 244       & ---       & 332       \\
Poland               & 200       & 192       & 363       & ---       & 494       \\
Portugal             & 39        & 38        & 100       & ---       & 136       \\
Romania              & 14        & 13        & 98        & ---       & 133       \\
Slovakia             & 18        & 17        & 56        & ---       & 76        \\
Slovenia             & 6         & 6         & 24        & ---       & 33        \\
Spain          & 389       & 373$^{c}$       & 502       & 375       & 683       \\
Sweden               & 282       & 270       & 139       & ---       & 189       \\
\midrule
\textbf{EU-27 total} & \textbf{4{,}481} & \textbf{4{,}298} & \textbf{4{,}298} & \textbf{1{,}891} & \textbf{5{,}847} \\
\bottomrule
\multicolumn{6}{p{15cm}}{\footnotesize
%Sources: \citet{iab_adex_2024} for national digital advertising market data.
%M1: SWD(2018) 81 EU-27 base (EUR 4{,}700\,M EU-28 $\times\,0.70$ excl. UK, per fn.~71) scaled to 2024 via IAB EU-27 advertising market growth 2022--2024 (factor 1.363, CAGR~$\approx$~17\%/yr, post-Covid recovery peak) \citep{swd_2018_81,iab_adex_2024}; allocated by IAB 2024 country shares.
%M2: $R_i = 3\% \times \bar{k} \times A_i$, $\bar{k} = 2.260$; 2024 base.
%M3: M2 EU-27 total reallocated by Eurostat isoc\_ec\_ib20 demand-side weights (individuals ordering online, 2024).
%M3~(2028): M3 figures projected to 2028 at a CAGR of 8\%/yr, derived from Statista Market Insights digital services market forecasts as reported in \citet{ep_budg_digital_levy_seminar_jun2026}; a baseline assumption; the 2022--2024 historical rate of approximately 17\%/yr reflected post-Covid recovery.
%Conf.: confirmed administrative revenues at national rates.
$^{a}$~Austria scope incompatible (advertising-only at 5\%).
$^{b}$~M1 and M2 not available: Cyprus, Luxembourg, Malta are not covered by IAB Europe AdEx Benchmark. M3 derived from Eurostat demand-side weights applied to the EU-27 aggregate; M1 and M2 EU-27 totals therefore reflect EU-24 coverage only.
$^{c}$~France and Spain are M2 anchors.
$^{d}$~Italy 2025 exceeds M2 due to threshold removal.}\\
\caption{Classic three-category DST revenue by method and member state (EUR m, 3\% rate)}\label{tab:compare}
\end{dsttab}

The M3~(2028) column projects each member-state estimate to the first full year of MFF~2028--2034 operation.
The projection applies a baseline CAGR of 8\%/yr to the M3~(2024) base, derived from Statista Market Insights forecasts for EU digital services markets~\citep{statista_digital_market_insights_2025}, against a historical 2022--2024 rate of approximately 17\%/yr that reflected post-Covid recovery.
At this rate the EU-27 total reaches EUR~5.8\,bn at 3\% in 2028. % AF: I'd remove "at 3\%"
The IAB EU-27 advertising market growth factor of 1.363 over 2022--2024 (CAGR~$\approx$~17\%/yr) used in M1 is retained as the calibration anchor for the 2024 base but is not extrapolated forward.

M2 and M3 diverge at some places.
M2 weights by IAB national advertising market shares: Germany, with the largest digital advertising industry in the EU, ranks first.
M3 weights by Eurostat online buyer counts: France, with the largest online buyer population, ranks first.
Poland and Romania rise substantially under M3; Sweden and Denmark fall.

The EU-27 aggregate is unaffected: both methods yield EUR~4.3\,bn.
What varies is the distribution across member states.

Both methods are proxies for the taxable base.
M2 follows IAB advertising market shares, a variable that reflects advertiser demand and correlates with user location but is also influenced by the size of each country's corporate advertising sector; Germany illustrates this, with a large advertising market that reflects B2B and industrial advertising intensity, not simply user volume.
M3 follows Eurostat online buyer counts, a consumer-side variable that more directly measures where users transact across the three taxable categories.
For an instrument whose nexus is user location, the consumer-side proxy is the more internally consistent choice.
Neither is a direct measure of taxable revenues; M3 is the best available public approximation.


\section{Effects and responses}

\subsection{Practical viability}
\label{sec:effects:viability}

Two questions bear on the practical viability of a DST: who ends up paying, and whether firms can reduce their exposure by restructuring.
Six years of national experience alongside recent academic work on platform taxation shed light on both.

National experience corroborates both.
In France, Amazon notified marketplace sellers of a 3~per cent fee increase on the day the TSN was enacted.
In Italy, Google applies a 2.5~per cent operational surcharge to advertising services~\citep{olmastroni_2025}.
The tax is not borne by the platforms; it is passed to the businesses that depend on them: marketplace sellers and advertisers, many of them SMEs, and from there in part to end consumers.
This is a genuine design constraint for a gross-revenue levy, and it characterises every national DST already in force, not only the proposed EU instrument.

Langenmayr and Muddasani~\citep{langenmayr_muddasani_2026}, using Amazon's Fulfilled by Amazon fee schedules across France, Spain, Italy, and the United Kingdom following each DST introduction show that platforms do not necessarily absorb the tax.
In France, Spain and the UK, between 80 to more than 100\%  was passed on in fee increases to third-party sellers.
Sellers then passed those higher costs on to consumers, with retail prices rising by more than the fee increase.
Italy is the exception: fee increases there were small and not statistically significant, showing that pass-through is a commercial choice rather than a structural necessity. % AF: chelou qu'il y ait pas les mêmes effets en Italie, on sait pourquoi ?

For the online advertising market specifically, Lassmann et al.~\citep{lassmann_2025} provide the closest empirical test.
They use Facebook's 2016 accounting restructuring as a natural experiment: where effective corporate tax rates rose, ad prices rose and advertisers bore the cost, consistent with a model predicting that the platform restricts ad supply in higher-taxed markets. % AF: est-ce que ce serait pas plus clair de dire que model's prediction c'est que platform ne réduit pas ses prix hors taxe / ses profits?
The instrument is a corporate income tax, not a DST, but the authors note the price response is consistent with the advertising surcharges platforms applied when national DSTs were introduced.


On avoidance, the scope for restructuring is structurally limited.
Firms can in principle reclassify revenue streams or adjust intragroup allocations, but a DST based on user location cannot be avoided by moving where revenues are booked: users cannot be relocated.
National experience confirms this: revenues in all four countries with data grew consistently after introduction, against expectations of significant avoidance~\citep{ep_budg_digital_levy_seminar_jun2026}.
Fears of systemic volume contraction have similarly not materialised; digital advertising spend grew at rates comparable to non-DST markets.

On revenue stability, a tax on revenues is less sensitive to the business cycle than a tax on profits.
Corporate income tax receipts can fall sharply in a downturn; digital activity and user engagement tend to hold up.
The consistency of DST receipts in France, Italy, Spain, and Austria since introduction supports this~\citep{ep_budg_digital_levy_seminar_jun2026}.
Spend growth does not resolve the question of impression volumes, however: higher CPMs with stable advertiser budgets mean fewer impressions per euro, a channel the available data leave unresolved. % AF: je comprends pas pq tu dis ça: pourquoi stable advertiser budgets alors que tu dis plus haut que les budgets augmentent suite aux taxes ?


\subsection{Fiscal contribution}
\label{sec:feasibility:revenue}

The EUR~4.3\,bn estimate at 3\% represents additionality: new public revenue from a sector that currently contributes far less in corporate tax than comparable traditional sectors. % AF: "at 3\%" is not grammatically correct
Germany, accounting for roughly 18\% of EU-27 digital advertising and currently operating no national DST, would for the first time see revenues from digital services in its territory contribute to the European budget.
The same applies to the Netherlands, Belgium, and Poland, all of which have substantial digital markets and no existing instrument capturing that base.

For member states that already operate national DSTs, replacement has a different value.
France's TSN collects EUR~756\,M per year while keeping France on the USTR Section~301 investigation list.
An EU instrument removes that bilateral exposure for all DST countries at once.
The fiscal transfer into the EU budget is offset by each state's share of the expenditure it finances; where the net balance proves burdensome, the MFF contribution framework offers the standard adjustment mechanism.
The risk reduction is immediate; the fiscal arithmetic is negotiable.

\subsection{A unified instrument}
\label{sec:feasibility:replace}

An EU instrument replacing rather than adding to the national DSTs eliminates the overlap problem. % AF: be more explicit than "overlap pb"
Member states that currently operate a DST phase it out as the EU instrument enters into force.
The compliance cost for in-scope companies falls: one rate, one consolidated EU revenue base, one reporting framework, instead of up to eight national regimes with divergent definitions and nexus rules.
Inconsistent national scope definitions create arbitrage that a uniform directive eliminates~\citep{eu_commission_proposal}.

The replacement logic is directly precedented by the OECD Pillar~1 standstill: countries committed to withdraw national DSTs upon Amount~A implementation.
That trigger was never pulled.
An EU instrument offers the same outcome at EU level, with the difference that the EU has jurisdiction to make the commitment binding. % AF: dernière phrase pas claire

\subsection{Collective leverage}
\label{sec:strategic}

An EU-level DST is not merely a technically superior version of the national instruments now in force.
The shift from individual member-state action to a collective EU instrument changes the strategic position of European governments vis-\`{a}-vis the United States in a way no individual state can replicate.

Section~301 of the Trade Act of 1974 works by threatening retaliatory tariffs on the DST country's goods exports~\citep{ustr_section301_dst,skadden_feb2025}.
France faces renewed Section~301 investigations since February 2025; Italy is similarly listed.
But Section~301 tariffs on French exports legally implicate the entire EU single market and trigger the EU's Trade Enforcement Regulation countermeasure framework.
An EU-level DST replaces seven individual bilateral targets with a single EU-level interlocutor backed by the weight of the single market.

The 2025 Section~899 episode confirms the leverage differential~\citep{crs_if13023}.
Section~899's threat alone, before enactment, produced the G7 ``shared understanding'' of 26 June 2025 and the subsequent OECD side-by-side package.
Section~301, relaunched in February 2025 and escalated to a 100\% tariff threat in June 2026, has produced no DST withdrawal in Europe.
Canada's repeal in June 2025 is the single exception: it was driven by direct bilateral negotiation pressure on a country without the EU's collective trade architecture~\citep{canada_finance_rescission_jun2025}.
The asymmetry is the argument.

The MFF window reinforces the strategic case.
Incorporating the DST into the own-resources package changes the political economy of blocking.
A country that benefits from EU structural and cohesion funding faces a different calculation when opposing a DST means opposing the entire MFF settlement.
Collective leverage at the budget table is a tool individual national DSTs do not carry.


\section{Feasibility}
\label{sec:feasibility}

\subsection{Legality}
\label{sec:feasibility:legal}

Article~113 TFEU provides the legal base for the harmonisation of indirect taxes.
A DST structured as a levy on revenues rather than profits falls within its scope.
The French Conseil constitutionnel confirmed on 12 September 2025 that user location is a constitutionally sufficient nexus for the TSN~\citep{conseil_constitutionnel_sept2025}.
A well-drafted EU directive replicating that architecture inherits the same legal characterisation.

The DST is not a tax on income within the meaning of bilateral tax treaties.
France and the UK have confirmed this status explicitly for their national instruments: foreign tax credit claims in the United States are not triggered, and treaty non-discrimination provisions do not apply.
These are the two key defences against the main legal challenges from treaty partners.

France, Germany, Italy, Spain, the Netherlands, Belgium, and Poland together represent the dominant share of EU GDP and of the EU digital revenue base; the political weight of this coalition creates strong pressure on non-participating states to align. % AF: why these countries?

\subsection{Member state support}
\label{sec:feasibility:governments}

The premise that a European DST faces unanimous member-state opposition is not accurate.
Several governments have explicitly conditioned their national positions on the absence of an EU or OECD solution.

The Netherlands stated in 2023 that ``an EU DST should be considered as an alternative to Pillar One, Amount A if a global agreement is not reached.'' That agreement has now collapsed.
Belgium's coalition agreement of January 2025 commits to a 3\% national DST by 2027 contingent on the absence of EU or OECD consensus, making the EU route the stated preference and the national instrument a fallback.
Poland's Finance Ministry proposal of January 2026 is framed explicitly as temporary pending an OECD solution; Slovakia's MIRDI proposal of August 2025 uses identical language.
In both cases the national proposal is a placeholder, not a policy preference.

The European Parliament has moved beyond endorsement to conditionality.
The interim MFF report adopted in April 2026~\citep{ep_interim_report_mff_2028} names a digital services levy as the first alternative own resource if the Council rejects the Commission's package, and Parliament has stated it will not consent to an MFF without sufficient new own resources.
Germany, the largest potential beneficiary of the instrument (a large digital market, no existing DST revenue to forgo), has not opposed an EU-level instrument.

\subsection{Public opinion backing}
\label{sec:feasibility:democratic}


Public opinion reinforces the procedural mandate.
A Flash Eurobarometer conducted across all 27 member states in April 2025 found that 80 per cent of EU respondents agreed that large multinational companies should be required to pay a minimum amount of tax in each country where they operate (44 per cent strongly, 36 per cent somewhat); the range across member states ran from 67 per cent to 87 per cent, with no country below a majority~\citep{flash_eb_taxation_2025}.
A Kieskompas poll conducted in November 2018, as reported by the S\&D Group in the European Parliament, found more than 80 per cent of respondents in France, Germany, the Netherlands, Denmark, Sweden, and Austria in favour of a fair tax on technology companies~\citep{kieskompas_sd_2018}.
A more recent EU-wide survey (Polling Europe Euroscope, July 2026, n=5{,}048) found 39 per cent in favour of an EU-level digital levy, 29 per cent preferring national-level action, and 11 per cent opposed to any such measure~\citep{polling_europe_euroscope_13_2026}; taken together, the two surveys show that more than two thirds of respondents supported taxing digital revenues in some form.
The Tax Justice Network~\citep{tax_justice_network_sotj2024} has called for stronger taxation of digital revenues. Oxfam~\citep{oxfam_dst_dec2018} criticised the Commission's 2018 proposal as insufficient, arguing it let tech giants off the hook and that more ambitious measures were needed.

\section{Principal objections answered}
\label{sec:feasibility:objections}

\subsection*{The US will retaliate.}
It already is.
Section~301 investigations are active against France, Austria, Italy, Spain, and the UK~\citep{skadden_feb2025}; the question is not whether to incur retaliation risk but how to manage it.

The 2021 termination of those actions without a single DST withdrawal~\citep{wilmerhale_nov2021} showed that the threat carries a political cost for the US as well: imposing 25\% duties on French wine and Italian handbags to discipline a tax on Google is disproportionate, and the Biden administration did not follow through.
The 2025 context is less favourable.
The Trump administration has demonstrated a higher tolerance for trade escalation, Section~301 has been relaunched with a 100\% tariff threat, and whether the same logic holds is genuinely uncertain.

What is not uncertain is that a fragmented set of national DSTs is weaker than a collective EU instrument in this environment.
Section~301 tariffs on any member state's goods raise the cost across the entire EU single market and activate the EU's Trade Enforcement Regulation countermeasure framework.
Canada repealed its DST in June 2025 under direct bilateral pressure; Canada does not have that architecture~\citep{canada_finance_rescission_jun2025}.
The higher the retaliation risk, the stronger the argument for consolidating seven bilateral exposures into one EU-level interlocutor.

\subsection*{The instrument will be successfully challenged as discriminatory.}
The objection is predictable; it has also been anticipated in the design.

A levy on gross revenues from specified services is not an income tax within the meaning of bilateral tax treaties.
Treaty non-discrimination provisions do not apply; foreign tax credit claims are not triggered.
France and the United Kingdom have both confirmed this characterisation before their respective courts, and a well-drafted EU directive inherits both precedents.

Taxing rights follow user location, a principle derived from the OECD's own user value creation framework under BEPS Action~1.
An instrument that taxes revenues generated by serving users in a territory, applied identically to any company serving those users, cannot easily be characterised as a measure directed at foreign companies.
No company has overturned this principle before any EU member state court in five years of operation.

Finally, the threshold of EUR~750~million in global revenues is a size criterion and not in itself discriminatory; European companies meeting it are in scope on the same terms.
Booking.com, Just Eat Takeaway, and Spotify approach or meet that level: their presence in scope is an inconvenience for some member states and a legal asset for the instrument.


\subsection*{Unanimity is impossible.}
Some member states have structural fiscal interests in blocking; Ireland, Luxembourg, Malta, and Cyprus are the principal cases, and Malta has flagged a veto explicitly~\citep{tax_foundation_dst_tracker_2026}.
The MFF package is the lever.
A country that benefits from structural and cohesion funding faces a different calculation when opposing the DST means opposing the entire budget settlement.
The Parliament's conditionality makes that logic compelling: the body with budgetary consent authority has said it will not approve an MFF without sufficient new own resources.

\subsection*{The DST taxes revenues rather than profits and may reduce investment.}
The distortion the DST corrects is larger than the one it introduces.
The 13.7 percentage-point effective tax rate gap between digital and traditional firms is a structural artefact of rules that predate the digital economy, not a market outcome.
A 3\% revenue levy that partially closes that gap is corrective rather than purely distortionary.
The gross revenue base compounds this: royalty flows to low-tax IP holding entities can eliminate taxable profit entirely, but they cannot reduce advertising revenues or intermediation commissions.
The instrument resists precisely the avoidance mechanism that makes corporate income tax ineffective against this sector.

The investment concern deserves a direct answer.
Revenue-based taxation creates a liability for firms that are loss-making or operating at thin margins.
The EUR~750 million global threshold screens out exactly those firms; the companies in scope are large, structurally profitable platforms.
Five years of national data show no suppression of digital investment or platform expansion in France, Italy, Spain, or Austria.

\subsection*{DST member states would transfer revenues from their budgets to Brussels.}
A member state operating a national DST does transfer that revenue stream to the EU budget.
The trade is more favourable than the headline revenue figure suggests.
Bilateral Section~301 exposure is replaced by collective EU-level interlocution: a tariff threat against any member state triggers the Union's Trade Enforcement Regulation countermeasure framework, which no individual country can invoke alone.
The compliance arbitrage that allows platforms to minimise liability by exploiting divergences between national regimes disappears under a uniform instrument.
EU enforcement capacity replaces national administration against counterparties that individual tax authorities cannot match.
The national instrument, always transitional by design, is replaced by a more durable one whose design the member state has a direct role in shaping. % AF: I thought it was also temporary
Should the fiscal transfer prove burdensome for particular economies, the MFF framework provides a collective adjustment lever.
In aggregate the accounting remains favourable: all member states share in the expenditure the instrument finances, and the net fiscal flow is positive for the Union as a whole.

% -------------------------------------------------------

% -------------------------------------------------------
\section{Extensions}
\label{sec:extensions}

Extensions serve a dual purpose: they raise EU revenues, and by demonstrating that the Union can deepen its digital tax architecture without a multilateral agreement, they create an incentive for trading partners to resume the negotiation process.
The three-category instrument represents a conservative starting point.
Two extensions merit consideration once the core instrument is established.

\textbf{Rate.} At 3\%, the three-category base yields approximately EUR~6\,bn by 2028. The instruments already in force include rates of 5\% (Austria), 6\% (France), and 7.5\% statutory (Hungary). At 5\%, the EU-27 yield scales proportionally to approximately EUR~10\,bn; at 6\%, matching the current French rate, to approximately EUR~12\,bn. % AF: ça représente quel % rapporté au profit?

\textbf{New categories.} Thomadakis~\citep{ep_budg_digital_levy_seminar_jun2026} identifies cloud computing services as a fourth major category of digital value extraction, one expected to grow at the highest rate of any digital revenue stream. The present design does not reach cloud infrastructure, software-as-a-service subscriptions, or AI inference APIs sold to EU-based customers. Extending to these services would follow the same EU-user nexus logic as the existing three categories and raise the tax base materially, while requiring careful legislative drafting to prevent double-counting with existing intermediation revenues.

% -------------------------------------------------------
\section{A Pigouvian advertising levy}
\label{sec:pigouvian}
\subsection{The external cost of digital advertisements}

We saw in Section~\ref{sec:effects:viability} that pass-through of DST to clients of big digital companies is important.
The tax is borne by downstream market participants, final clients and companies paying for ads.
Platform profits remain unaffected.
This partial failure, however, opens a different reading of the instrument's value.
The economic impact falls partly on digital advertising, a sector with documented externalities.
This is exactly the economic impact on this sector that makes this instrument valuable.

Criticism of advertising is not necessarily new, % AF: ici tu peux citer Dixit, Avinash, and Victor Norman. 1978. “Advertising and Welfare.” The Bell Journal of Economics 9 (1): 1–17. https://doi.org/10.2307/3003609. L'argument c'est que la captation de l'attention est un jeu à somme nulle qui mène à un excès de pub
but as the sector is booming, an entire attention economy has been developed.
More and more energy is put to keep people captive on platforms and to present them ads skillfully engineered for them.
Artificial Intelligence will make this even easier to implement and will amplify such externalities in a hardly predictable manner.
We describe three categories of documented externalities, following Acemoglu and Johnson~\citep{acemoglu_johnson_2024_ad_tax}.

Platforms optimize for engagement through emotionally charged content.
Outrage and indignation are the most effective levers to generate strong engagement and condition users to engage in purchases with a higher success rate.
This has documented effects, particularly among young people, in terms of mental health, addiction, body image, self harm, eating disorders, sleep disturbances and bullying.
Users bear these harms without viable exit: network effects bind them to platforms where their social life is.
Platforms choose network architectures showing content only to users who are predisposed to agree with it.
To attain maximum ideological homophily, algorithms are knowingly built to allow unreliable content to spread seamlessly within a community of users who share the same biases~\citep{acemoglu_ozdaglar_siderius_2024}.
Polarizing and violent discourse spread as a product of this optimisation, with documented effects in terms of radicalization and misinformation.
 

Consumers must be profiled so that ads fitting their characteristics can be presented to them.
The distribution algorithms are fed by data also obtained by monitoring user behaviour.
Advertisers pay for access to users' attention.
Users pay with their time, their cognitive engagement, and their data.
Acemoglu and Johnson~\citep{acemoglu_johnson_2024_ad_tax} argue that the advertising revenue ecosystem encourages valuable engineering and artificial intelligence talent to concentrate on surveillance, engagement optimization, and the automation of human interactions.
Such skills could contribute more to global well-being if they were directed toward supporting workers and improving productivity.


Digital advertisements also distort market structure.
Targeted advertising strongly influences consumers’ willingness to pay by leading them to overestimate the quality of a particular product. 
Even if not all consumers are affected in the same way, the fact that some are weakens price competition: margins and prices for artificially differentiated products rise. 
Competition between platforms should, in theory, drive prices down and improve user well-being, but this is not the case. 
All this advertising engineering promotes products regardless of their quality. 
Economic actors are not encouraged to improve their products but rather to multiply the signals directed at unsuspecting users with a high ad load.
Competition does not incentivize a reduction in this manipulation: each platform can maintain a high ad load without losing its unsuspecting users to a competitor.
The equilibrium remains separating, with low welfare, even with multiple competing platforms~\citep{acemoglu_huttenlocher_ozdaglar_siderius_2024}.


Acemoglu and Johnson propose a rate of 50 per cent to force a structural switch away from advertising-dependent business models entirely.
Their argument is not derived from a precise measurement of marginal external cost.
Like tobacco excise taxes, the rate is calibrated to change behaviour: the harm is documented and large and the market will not self-correct.
A price signal strong enough to alter incentives is the appropriate instrument.

National DSTs already in place at 3 per cent were not designed to internalise negative externalities.
On the contrary, little market change was observed.
At 15 per cent, the relative attractiveness of subscription-based alternatives shifts enough to matter.
We propose to start at 15 per cent to open at a politically tractable entry point with a clear signal that the trajectory is escalating.
The rate should rise as the instrument demonstrates its effects, the evidentiary case matures, and the political conditions for a more ambitious target develop.
This is more an incentive to build different business models than a revenue instrument.
Nonetheless, this section develops that instrument as a standalone EU own resource.
The levy presented here is independent of the three-category DST presented above, a separate option for a Union with the ambition to steer its economy toward collective welfare.
The advertising revenue base of the Pigouvian levy overlaps with the first category of the DST; the two instruments are best understood as sequential options, and their simultaneous application would bring the combined levy on digital advertising to 18 per cent, 
a question that would need to be addressed explicitly in the relevant legislation. % AF: I'd remove the last line, or replace it by "a rate close to the Pigouvian-only tax".

\subsection{A digital ad tax at 15\,\%}

The EU-27 digital advertising market reached EUR~63.4\,bn in 2024, measured by the IAB Europe AdEx Benchmark~\citep{iab_adex_2024}.
The market is highly concentrated: Google (Search and YouTube), Meta, and Amazon together account for approximately 75\,\% of revenues.
We propose a threshold of EUR~100\,M in annual EU advertising revenues, derived from the USD~500\,M threshold proposed by Acemoglu and Johnson~\citep{acemoglu_johnson_2024_ad_tax}, % AF: c'était pas à l'échelle mondiale eux ?
scaled to the EU digital advertising market, which represents roughly one quarter of the United States market by revenue.
This excludes the long tail of smaller publishers and ad-tech intermediaries and reduces administrative burden without protecting the platforms whose business models drive the documented harms.
A revenue-from-advertising nexus, rather than a global turnover threshold, also makes the instrument harder to characterise as targeting a specific set of US firms: any platform generating EUR~100\,M in EU advertising revenues falls within scope regardless of corporate origin.

The pass-through analysis in Section~\ref{sec:effects:viability} was based on data from the digital advertising sector.
For revenue estimation, full pass-through is assumed, though the actual incidence may be lower or may involve over-shifting.
With this assumption, a 15\,\% revenue tax raises the advertiser price by $\tau/(1-\tau) \approx 17.6\,\%$.
Under a standard constant-elasticity demand function $Q \propto P^{\varepsilon}$, net tax revenue becomes:
\[
T = \tau \cdot B_0 \cdot \left(\frac{1}{1-\tau}\right)^{\!1+\varepsilon}
\]
where $B_0$ is the pre-tax advertising base.

No peer-reviewed study directly estimates the price elasticity of advertiser demand in response to a gross-revenue tax on EU digital advertising.
\citet{pellefigue_2019} notes that advertising demand is elastic without providing a numerical estimate.
\citet{dandria_2023} derives conditions under which an ad valorem digital advertising tax can paradoxically raise equilibrium ad volume, depending on demand curvature, but does not estimate $\varepsilon$ empirically.
Large-scale platform experiments~\citep{blake_nosko_tadelis_2015} suggest demand for non-captive ad inventory is meaningfully elastic; short-run budget rigidity makes it less so.
The range $\varepsilon \in [-0.5,\,-1.5]$ spans what the available literature implies; Table~\ref{tab:revenue_pig} presents the resulting yields. % AF: ça sort d'où cette fourchette ? Tu devrais aussi le calculer pour un taux de 50% vu que c'est l'objectif ultime. TODO
\begin{dsttab}{@{}lrrr@{}}
\toprule
Scenario & $\varepsilon$ & Net yield (EUR~bn) & Market contraction (EUR~bn) \\
\midrule
\endfirsthead
\multicolumn{4}{l}{\textit{(continued)}} \\
\toprule
Scenario & $\varepsilon$ & Net yield (EUR~bn) & Market contraction (EUR~bn) \\
\midrule
\endhead
\endfoot
\endlastfoot
Low response  & $-0.5$ & 14.0 &  6.7 \\
Central       & $-1.0$ & 12.9 & 12.9 \\
High response & $-1.5$ & 11.9 & 18.7 \\
\bottomrule
%\multicolumn{4}{p{13cm}}{\footnotesize Zero-elasticity yield: EUR~15.2\,bn. Central scenario ($\varepsilon = -1.0$): EUR~12.9\,bn.
%Base: IAB Europe AdEx Benchmark \citeyear{iab_adex_2024} 2024 market size of EUR~63.4\,bn projected to 2028 at a CAGR of 8\,\%/yr (EUR~86.2\,bn), consistent with Statista Market Insights forecasts for EU digital services markets~\citep{statista_digital_market_insights_2025}.
%Full pass-through assumed.
%Market contraction: reduction in advertiser spending at pre-tax prices, $\Delta B = B_0\,(1-(1-\tau)^{-\varepsilon})$.} \\
\caption{Projected net yield and market contraction at 15\,\% on EU-27 digital advertising revenues (2028)}\label{tab:revenue_pig}
\end{dsttab}

The revenue shortfall relative to the static baseline is the mechanism working, not a cost to minimise: a contraction in advertising spend is fewer resources directed toward attention maximisation, and the fiscal yield of EUR~12 to 14\,bn per year at 2028 market size is an appreciable by-product of the correction.

Table~\ref{tab:country_pig} distributes the central-scenario yield across EU member states using the 2024 IAB Europe country figure~\citep{iab_adex_2024} projected to 2028 at an 8\,\% annual CAGR.
This is then multiplied by $\tau$ as the formula in Table~\ref{tab:revenue_pig} simplifies in the case $\varepsilon = -1.0$.

\begin{dsttab}{@{}p{4cm}r@{}}
\toprule
\textbf{Member state} & \textbf{Yield (EUR\,m)} \\
\midrule
\endfirsthead
\multicolumn{2}{l}{\textit{(continued)}} \\
\toprule
\textbf{Member state} & \textbf{Yield (EUR\,m)} \\
\midrule
\endhead
\midrule
\multicolumn{2}{r}{\textit{(continued on next page)}}
\endfoot
\endlastfoot
Germany        & 3{,}652 \\
France         & 2{,}287 \\
Italy          & 1{,}205 \\
Spain          & 1{,}122 \\
Netherlands    &   882 \\
Sweden         &   813 \\
Poland         &   577 \\
Austria        &   446 \\
Czech Republic &   355 \\
Denmark        &   350 \\
Belgium        &   248 \\
Ireland        &   216 \\
Finland        &   181 \\
Hungary        &   153 \\
Portugal       &   114 \\
Greece         &    74 \\
Slovakia       &    51 \\
Romania        &    40 \\
Croatia        &    36 \\
Lithuania      &    32 \\
Estonia        &    29 \\
Latvia         &    27 \\
Bulgaria       &    24 \\
Slovenia       &    18 \\
\midrule
\textbf{EU-27} & \textbf{12{,}930} \\
\bottomrule
%\multicolumn{2}{p{9cm}}{\footnotesize Central scenario ($\varepsilon = -1.0$): yield computed as $T = \tau B_0$ where $B_0$ is the 2024 IAB Europe country figure~\citep{iab_adex_2024} projected to 2028 at an 8\,\% annual CAGR (four-year factor: 1.361). At $\varepsilon = -1.0$, the formula in Table~\ref{tab:revenue_pig} simplifies to $T = \tau B_0$. Ireland's figure reflects the user-location nexus; revenue booked by EU subsidiaries of US platforms is allocated to the country of the user.} \\
\caption{Pigouvian advertising levy: per-member-state yield at central scenario, 2028 (EUR million)}\label{tab:country_pig}
\end{dsttab}

\subsection{International positioning}

A Pigouvian advertising levy operates on entirely different legal ground from the OECD framework.
It targets externalities generated by the attention-capture business model, not the profit-shifting that Pillar Two addresses; the two instruments are complementary and do not interact.
It is an excise on an externality-generating activity, not a digital services tax, and Pillar One repeal commitments do not apply.
No multilateral negotiation is required or sought: the externality rationale is self-standing, and the EU is not conditioning the instrument on an international process.

Governments were pressed under Section~301 in 2025 because their instruments fell predominantly on a small set of US-headquartered platforms~\citep{ustr_section301_dst}; a Pigouvian advertising tax grounded in documented externalities, with a threshold of EUR~100\,M set by market concentration rather than national origin, does not meet that discrimination standard.
Section~899 remains a threatened fallback following its removal from the One Big Beautiful Bill Act; a genuine externality rationale makes the instrument harder to classify as an ``unfair foreign tax,'' and the EU's collective legal standing under the Trade Enforcement Regulation raises the cost of any retaliation substantially above what individual member states face~\citep{skadden_feb2025}.
A threshold defined by size rather than national origin and a legislative preamble documenting the externality rationale are the design choices that matter for both defences.

The instrument does not require international coordination to operate.
Acemoglu and Johnson~\citep{acemoglu_johnson_2024_ad_tax} nonetheless call for G7 harmonisation of such a levy, on the grounds that single-jurisdiction action leaves platforms room to shift activity.
An EU-level introduction would constitute the European contribution to that process, and evidence that a tax instrument can redirect the digital economy away from attention extraction and toward applications that genuinely serve users.

% -------------------------------------------------------
\section*{Conclusion}
\addcontentsline{toc}{section}{Conclusion}


Six years of international efforts have yielded no results on the issue of taxing digital companies multilaterally, leading the landscape of digital services taxes to take shape around national initiatives. 
This currently results in coverage gaps across the European single market. 
We propose for the EU to establish a digital services tax on revenues. 
Although imperfect in many respects, it is designed to be repealed as soon as international regulations have been adopted and implemented. 
With its DST, the EU is not simply waiting passively for better legislation; it is calling on other countries to return to the negotiating table regarding an agreement on profit taxes.

We also propose a tax on digital advertising, justified by the sector’s negative externalities.
The explicit goal is to shift the dynamics of this economic sector towards a direction more conducive to the common good.

The political challenges are real: between diverging structural interests and announced vetoes, unanimity appears difficult to achieve.
This is not unique to the instruments proposed here: many proposals for new EU own resources have supporters and detractors among Member States.
That is why the constraints and stakes of each party must be addressed collectively.
The MFF negotiations appear to be the forum where these considerations can be weighed and negotiated as a whole across proposals.
The two options presented here feed into this discussion.

These proposals go beyond a simple budgetary plan by contributing to European unity, building a global positioning at the international level, and steering the Union’s economic policy in a direction that reflects collective priorities.
From a budgetary perspective, both of our options will contribute to the repayment of NextGenerationEU: approximately EUR~6\,bn for the first, and EUR~13\,bn for the second.

\newpage
\bibliographystyle{apalike}
\bibliography{DST}

\end{document}
